% =========================================================================
% Capítulo 4: Cuerpo rígido
% Tema 01: Cuerpo rígido
% Problema 15: Placa rectangular en una varilla apoyada en rodamientos
% Origen: Ayudantía final (Última ayudantía), 2026 (problema MIT)
% =========================================================================
\subsubsection{Placa rectangular en una varilla apoyada en rodamientos}

Una placa rectangular de masa $M$, longitud $b$ y altura $a$ está fijada a una varilla sin masa
a través de su centro de masa y las esquinas de la placa, como se muestra en la Figura
\ref{fig:c4_e01_p15_placa_diagonal}. La varilla está apoyada en dos puntos separados por una
distancia $L$, como se muestra. La placa rota junto con la varilla con una velocidad angular
$\omega$.

% --- FIGURA PRINCIPAL -------------------------------------------------------
\begin{figure}[!ht]
    \centering
    \begin{tikzpicture}[scale=0.95]
        % --- COLORES ---
        \colorlet{bodycolor}{cyan!15}
        \colorlet{shaftcolor}{gray!50}
        \colorlet{bearinggray}{gray!30}
        \colorlet{housingcolor}{black!70}

        % --- RODAMIENTO EN CORTE (pedestal) ---
        \tikzset{
          bearing/.pic={
            % Base y soporte inferior (achurado)
            \draw[thick, fill=white] (-0.8,-3) -- (0.8,-3) -- (0.5,-0.7) -- (-0.5,-0.7) -- cycle;
            \fill[pattern=north east lines, pattern color=housingcolor] (-0.8,-3) -- (0.8,-3) -- (0.5,-0.7) -- (-0.5,-0.7) -- cycle;
            % Soporte superior (achurado)
            \draw[thick, fill=white] (-0.5,0.7) rectangle (0.5,1.2);
            \fill[pattern=north east lines, pattern color=housingcolor] (-0.5,0.7) rectangle (0.5,1.2);
            % Anillos exteriores
            \draw[thick, fill=bearinggray] (-0.3,0.5) rectangle (0.3,0.7);
            \draw[thick, fill=bearinggray] (-0.3,-0.5) rectangle (0.3,-0.7);
            % Anillos interiores
            \draw[thick, fill=bearinggray!50] (-0.3,0.2) rectangle (0.3,0.3);
            \draw[thick, fill=bearinggray!50] (-0.3,-0.2) rectangle (0.3,-0.3);
            % Elementos rodantes
            \draw[thick, fill=white] (0, 0.4) circle (0.1);
            \draw[thick, fill=white] (0,-0.4) circle (0.1);
            % Eje vertical central
            \draw[dash dot, thin, gray!80] (0, 1.5) -- (0, -3.2);
          }
        }

        % --- SUELO ---
        \fill[pattern=north east lines, pattern color=gray] (-4.5, -2.5) rectangle (5.5, -3);
        \draw[thick] (-4.5, -2.5) -- (5.5, -2.5);

        % --- RODAMIENTOS ---
        \pic at (-2.5, 0) {bearing};
        \pic at (2.5, 0) {bearing};

        % --- EJE DE ROTACIÓN (varilla) ---
        \draw[thick, fill=shaftcolor] (-3.5, -0.2) rectangle (4.5, 0.2);

        % --- PLACA (cuerpo rígido M) ---
        \begin{scope}[rotate=33.5]
            \draw[thick, fill=bodycolor, fill opacity=0.9] (-1.8, -1.2) rectangle (1.8, 1.2);
            \node at (-1.2, -0.7) {\Large $\mathbf{M}$};
        \end{scope}

        % --- Detalles del eje pasando a través de la masa ---
        \draw[thick, dashed, opacity=0.6] (-1.7, 0.2) -- (1.7, 0.2);
        \draw[thick, dashed, opacity=0.6] (-1.7, -0.2) -- (1.7, -0.2);

        % --- EJE GEOMÉTRICO DE ROTACIÓN ---
        \draw[dash dot, thick] (-4.5, 0) -- (6, 0);

        % --- SISTEMA DE COORDENADAS ---
        \draw[thick, -{Stealth[length=3mm]}] (0,0) -- (0, 3.5) node[right] {\Large $\mathbf{y}$};
        \draw[thick, -{Stealth[length=3mm]}] (4.5, 0) -- (5.5, 0) node[below right] {\Large $\mathbf{x}$};

        % --- VELOCIDAD ANGULAR ---
        \draw[thick, dashed] (4.8, 0.8) arc (90:270:0.3 and 0.8);
        \draw[thick, -{Stealth[length=3mm]}] (4.8, -0.8) arc (-90:90:0.3 and 0.8);
        \node at (5.3, 0.7) {\Large $\vec{\omega}$};

        % --- COTAS Y ETIQUETAS ---
        \draw[dashed, thick, gray] (-2.5, -3) -- (-2.5, -4.2);
        \draw[dashed, thick, gray] (2.5, -3) -- (2.5, -4.2);
        \draw[{Stealth[length=2.5mm]}-{Stealth[length=2.5mm]}, thick] (-2.5, -3.9) -- (2.5, -3.9) node[midway, fill=white] {\Large $\mathbf{L}$};
        \node at (-1.5, -1.2) {\Large $\mathbf{a}$};
        \node at (1.1, -1.2) {\Large $\mathbf{b}$};
    \end{tikzpicture}
    \caption{Placa rectangular montada en su diagonal (problema MIT).}
    \label{fig:c4_e01_p15_placa_diagonal}
\end{figure}

\begin{enumerate}
  \item ¿Es el eje de rotación un eje principal?
  \item ¿Cuál es el momento angular de la placa? Haz un boceto del vector momento angular en relación con la posición de la placa.
  \item ¿Qué fuerzas deben ejercer los soportes? Da una fórmula para las fuerzas en los dos soportes en función del tiempo.
  \item Si la varilla no careciera de masa, sino que tuviera una masa $M_r$, un radio $R_0$ y una longitud $L_2$, ¿cómo cambiaría tu respuesta?
\end{enumerate}

% --- SOLUCIÓN ---------------------------------------------------------------
\begin{solucion}{Placa montada en su diagonal}

\textbf{Planteamiento}\\
Ejes principales de la placa en su centro de masa $G$ (placa delgada):
\[
\vec{e}_1 \parallel b:\ I_1 = \tfrac{1}{12}Ma^2,\qquad
\vec{e}_2 \parallel a:\ I_2 = \tfrac{1}{12}Mb^2,\qquad
\vec{e}_3 \perp \text{placa}:\ I_3 = \tfrac{1}{12}M(a^2+b^2).
\]
La varilla va por la diagonal. Sea $d=\sqrt{a^2+b^2}$ y $\alpha$ el ángulo entre la diagonal y el lado $b$: $\cos\alpha=b/d$, $\sin\alpha=a/d$. Usamos una base que gira con la placa: $\vec{e}_x$ a lo largo de la varilla (fija en el espacio), $\vec{e}_s$ en el plano de la placa y perpendicular a la varilla, y $\vec{e}_n=\vec{e}_x\times\vec{e}_s$ normal a la placa. Según la figura (lado $b$ inclinado hacia arriba a la derecha):
\[
\vec{e}_x = \frac{b\,\vec{e}_1 - a\,\vec{e}_2}{d},\qquad
\vec{e}_s = \frac{a\,\vec{e}_1 + b\,\vec{e}_2}{d}.
\]

\textbf{a) ¿Es la varilla un eje principal?}\\
El tensor de inercia en la base $(\vec{e}_x,\vec{e}_s,\vec{e}_n)$ es
\[
[I_G] = \frac{M}{12d^2}
\begin{pmatrix}
2a^2b^2 & ab(a^2-b^2) & 0\\
ab(a^2-b^2) & a^4+b^4 & 0\\
0 & 0 & d^4
\end{pmatrix}.
\]
El término cruzado $I_{xs}=\dfrac{Mab(a^2-b^2)}{12(a^2+b^2)}$ es distinto de cero salvo que $a=b$.
\respuesta{No. La diagonal sólo es eje principal si la placa es cuadrada ($a=b$). En un rectángulo los ejes principales en el plano son paralelos a los lados.}
\patotip[squareglasses,shovel]{La diagonal no es eje de simetría del rectángulo; reflejar la placa respecto de la diagonal no la deja igual.}

\textbf{b) Momento angular}\\
Con $\vec{\omega} = \omega\vec{e}_x$: $\omega_1 = \omega b/d$, $\omega_2 = -\omega a/d$, $\omega_3=0$. Entonces
\[
\vec{H}_G = I_1\omega_1\vec{e}_1 + I_2\omega_2\vec{e}_2
= \frac{M\omega ab}{12d}\,\big(a\,\vec{e}_1 - b\,\vec{e}_2\big).
\]
En la base de la varilla:
\respuesta{$\displaystyle
\vec{H}_G = \underbrace{\frac{M a^2 b^2}{6(a^2+b^2)}\,\omega}_{H_\parallel}\vec{e}_x
\;-\;\underbrace{\frac{M ab\,(b^2-a^2)}{12(a^2+b^2)}\,\omega}_{H_\perp}\vec{e}_s,
\qquad |\vec{H}_G| = \frac{M ab\,\omega}{12}
$}
\begin{itemize}[nosep]
  \item $\vec{H}_G$ está \textbf{en el plano de la placa}, no sobre la varilla.
  \item Forma un ángulo $\beta$ con la varilla: $\cos\beta = \dfrac{2ab}{a^2+b^2}$, es decir $\beta = \tfrac{\pi}{2}-2\alpha$ (cero si $a=b$).
  \item Para $b>a$ se inclina hacia el lado opuesto al que se inclina el lado largo: tiende a quedar perpendicular a la dimensión larga, como en una barra inclinada.
  \item Está fijo a la placa, así que \textbf{gira con ella}: describe un cono de semiángulo $\beta$ alrededor de la varilla.
\end{itemize}

\begin{center}
\begin{tikzpicture}[>=Latex, scale=1.0]
  \def\b{3.6}\def\a{2.4}
  \begin{scope}[rotate=33.69]
    \draw[thick, fill=cyan!10] (-\b/2,-\a/2) rectangle (\b/2,\a/2);
  \end{scope}
  \draw[dash dot, thick] (-3.6,0) -- (3.8,0);
  \draw[->, very thick, azul] (0,0) -- (2.6,0) node[above, xshift=-2pt] {$\vec{\omega}$};
  % H: dirección 2ab x - (b^2-a^2) s  => ángulo atan(-7.2/17.28) = -22.6°
  \draw[->, very thick, myred] (0,0) -- (-22.6:2.9) node[right] {$\vec{H}_G$};
  \draw[gray, ->] (1.5,0) arc (0:-22.6:1.5) node[midway, right] {$\beta$};
  \draw[->, gray] (0,0) -- (0,1.2) node[left] {$\vec{e}_s$};
  \fill (0,0) circle (1.3pt) node[above left] {$G$};
  % fuerzas en los descansos
  \draw[->, thick, myred] (-3.0,-1.3) -- (-3.0,-0.05) node[pos=0, below] {$\vec{F}_A$};
  \draw[->, thick, myred] (3.3,1.3) -- (3.3,0.05) node[pos=0, above] {$\vec{F}_B$};
  \node[font=\small, align=center] at (0,-2.4) {Vista normal a la placa en $t=0$ ($\vec{e}_s=\vec{e}_y$).\\
  $\dot{\vec{H}}_G$ apunta hacia dentro del papel.};
\end{tikzpicture}
\end{center}

\textbf{c) Fuerzas en los soportes}\\
\emph{Fuerza neta.} $G$ está sobre el eje y no se mueve: $\sum\vec{F} = \vec{0}$. Con gravedad, cada soporte sólo carga $Mg/2$ (placa centrada entre soportes, como en la figura).

\emph{Momento.} $\vec{e}_x$ es fijo y $\vec{e}_s$ gira: $\dot{\vec{e}}_s = \vec{\omega}\times\vec{e}_s = \omega\,\vec{e}_n$. Con $\omega$ constante:
\[
\vec{\tau}_G = \dot{\vec{H}}_G = \vec{\omega}\times\vec{H}_G = -H_\perp\,\omega\,\vec{e}_n
= -\frac{M ab\,(b^2-a^2)}{12(a^2+b^2)}\,\omega^2\,\vec{e}_n.
\]
El momento es \textbf{normal a la placa} y gira con ella. Lo aplican los soportes con un par: $\vec{F}_A=-\vec{F}_B$, ambos perpendiculares a la varilla. Con $A$ en $x=-L/2$ y $B$ en $x=+L/2$:
\[
\vec{\tau}_G = \tfrac L2\vec{e}_x\times\vec{F}_B - \tfrac L2\vec{e}_x\times\vec{F}_A = L\,\vec{e}_x\times\vec{F}_B
\;\Rightarrow\;
\vec{F}_B = -\frac{H_\perp\omega}{L}\,\vec{e}_s.
\]
Las fuerzas dinámicas están \textbf{en el plano de la placa}. Tomando $t=0$ en la posición de la figura (placa en el plano $xy$, $\vec{e}_z$ saliendo del papel), $\vec{e}_s(t)=\cos\omega t\,\vec{e}_y+\sin\omega t\,\vec{e}_z$:
\respuesta{$\displaystyle
\vec{F}_A(t) = \frac{Mg}{2}\vec{e}_y + F_0\big(\cos\omega t\,\vec{e}_y+\sin\omega t\,\vec{e}_z\big),\qquad
\vec{F}_B(t) = \frac{Mg}{2}\vec{e}_y - F_0\big(\cos\omega t\,\vec{e}_y+\sin\omega t\,\vec{e}_z\big)
$\\[2mm]
$\displaystyle F_0 = \frac{M ab\,(b^2-a^2)}{12\,L\,(a^2+b^2)}\,\omega^2$}
\emph{Interpretación física.} Las ``fuerzas centrífugas'' sobre las dos mitades de la placa forman un par que tiende a girar la placa para poner su lado largo perpendicular al eje. Los soportes lo contrarrestan: el derecho empuja hacia $-\vec{e}_s$ y el izquierdo hacia $+\vec{e}_s$. El par dinámico es de magnitud constante pero gira con $\omega$: es una carga alternante sobre los rodamientos, y por eso se equilibran dinámicamente los rotores.

\patotip[glasses,tie=blue!70!black]{\textbf{Regla de oro:} si $\vec{H}$ no es paralelo a $\vec{\omega}$, los rodamientos lo pagan en cada vuelta. Por eso se balancean dinámicamente los rotores.}

\textbf{d) Varilla con masa $M_r$, radio $R_0$, longitud $L_2$}\\
La varilla es un cilindro coaxial con $\vec{\omega}$: su eje \emph{sí} es principal y su centro de masa está sobre el eje. Por lo tanto:
\begin{itemize}[nosep]
  \item Agrega una componente axial al momento angular: $\vec{H} = \vec{H}_G^{\text{placa}} + \tfrac12 M_r R_0^2\,\omega\,\vec{e}_x$. Es paralela a $\vec{\omega}$ y constante, así que no aporta a $\dot{\vec{H}}$.
  \item \textbf{Las fuerzas dinámicas en los soportes no cambian} ($F_0$ es la misma). $L_2$ y $R_0$ no aparecen en ellas.
  \item Sólo cambia la parte estática: cada soporte carga $(M+M_r)g/2$ si la varilla está centrada. Si no lo está, el reparto del peso depende de dónde quede su centro de masa respecto de los soportes.
  \item $I_{xx}$ total crece en $\tfrac12 M_rR_0^2$; esto sólo importaría si $\omega$ variara (par motor $=I_{xx}\dot\omega$).
\end{itemize}

\end{solucion}
